\documentclass[11pt]{article}

\usepackage[a4paper,margin=2.5cm]{geometry}
\usepackage{amsmath,amssymb,amsthm}
\usepackage{hyperref}

\title{Geometrische Strukturen in Einbettungen von Primzahlen in oktonionische Gitter}
\author{Thomas Hoffbauer}
\date{}

\newtheorem{definition}{Definition}
\newtheorem{proposition}{Proposition}
\newtheorem{lemma}{Lemma}
\newtheorem{remark}{Bemerkung}
\newtheorem{conjecture}{Vermutung}

\begin{document}
	
	\maketitle
	
	\begin{abstract}
		Wir untersuchen eine Einbettung der Primzahlen in ein diskretes Gitter der Oktonionen. Ziel ist es, additive Primzahlkonstellationen durch algebraische Inzidenzrelationen zu modellieren. Wir formulieren eine Assoziator-basierte Kopplungsbedingung und diskutieren deren strukturelle Eigenschaften sowie erste heuristische Konsequenzen.
	\end{abstract}
	
	\section{Einleitung}
	
	Die Verteilung der Primzahlen zeigt sowohl multiplikative als auch additive Struktur. Während klassische Methoden additive Eigenschaften analysieren, untersuchen wir hier eine Einbettung in eine nichtassoziative Algebra.
	
	\section{Oktonionischer Rahmen}
	
	Sei $\mathbb{O}$ die Algebra der Oktonionen mit Basis $\{e_0,\dots,e_7\}$.
	
	\begin{definition}
		Die Norm eines Elements $x=\sum a_i e_i$ ist definiert als
		\[
		N(x) = \sum_{i=0}^7 a_i^2.
		\]
	\end{definition}
	
	\begin{definition}
		Sei $\Lambda \subset \mathbb{O}$ ein diskretes Gitter (z.\,B. ganzzahlige Oktonionen). Für $n\in\mathbb{N}$ definieren wir
		\[
		S_n := \{x\in\Lambda : N(x)=n\}.
		\]
	\end{definition}
	
	\section{Einbettung der Primzahlen}
	
	\begin{definition}
		Für eine Primzahl $p$ definieren wir die Einbettung als
		\[
		\Phi(p) := S_p.
		\]
	\end{definition}
	
	\begin{remark}
		Diese Abbildung ist mengenwertig und vermeidet eine willkürliche Auswahl einzelner Repräsentanten.
	\end{remark}
	
	\section{Inzidenzstruktur}
	
	\begin{definition}
		Für $x,y,z \in \mathbb{O}$ definieren wir den Assoziator
		\[
		[x,y,z] := (xy)z - x(yz).
		\]
	\end{definition}
	
	\begin{definition}
		Für $x\in S_p$, $y\in S_q$ definieren wir die Assoziatorenergie
		\[
		\mathcal{A}(x,y) := \sum_{k=1}^7 \|[x,y,e_k]\|^2.
		\]
	\end{definition}
	
	\begin{definition}
		Ein Paar $(p,q)$ heißt geometrisch gekoppelt, wenn
		\[
		\exists x\in S_p, y\in S_q: \mathcal{A}(x,y)=0.
		\]
	\end{definition}
	
	\section{Grundlegende Struktur}
	
	\begin{lemma}
		Für $x\in S_p$, $y\in S_q$ gilt
		\[
		\|x-y\|^2 = p + q - 2\langle x,y \rangle.
		\]
	\end{lemma}
	
	\begin{proof}
		Folgt unmittelbar aus der Definition der Norm.
	\end{proof}
	
	\begin{proposition}[Heuristisch]
		Für kleine Primzahlen korrelieren kleine Differenzen $|p-q|$ mit kleinen Werten von $\mathcal{A}(x,y)$.
	\end{proposition}
	
	\section{Modellproblem}
	
	\begin{definition}
		Für $d\in\mathbb{N}$ sei
		\[
		\mathcal{C}_d := \{(p,q)\in\mathcal{P}^2 : q-p=d \text{ und geometrisch gekoppelt}\}.
		\]
	\end{definition}
	
	\section{Dynamische Heuristik}
	
	Wir betrachten den Raum
	\[
	X = (\mathbb{O}\otimes\mathbb{R})/\Lambda.
	\]
	
	Primzahlen definieren diskrete Abtastpunkte entlang eines Flusses.
	
	\begin{remark}
		Unter Annahme von Ergodizität würde eine positive Maßmenge stabiler Konfigurationen unendlich oft erreicht werden.
	\end{remark}
	
	\section{Vermutung}
	
	\begin{conjecture}
		Die Menge $\mathcal{C}_2$ ist unendlich.
	\end{conjecture}
	
	\section{Offene Probleme}
	
	\begin{itemize}
		\item Existenz einer kanonischen Auswahl $x_p\in S_p$
		\item Gleichverteilung der Punkte in $S_p$
		\item Nichttrivialität der Inzidenzbedingung
	\end{itemize}
	
	\section{Literatur}
	
	\begin{itemize}
		\item J. Baez: The Octonions
		\item Green--Tao: Arithmetic progressions in primes
		\item Hardy--Littlewood: Circle method
	\end{itemize}
	
\end{document}